% ==============================================================
%  informe.tex — Plantilla de informe para Labo 6
%  Compilar con: pdflatex informe.tex (+ biber si usás biblatex)
% ==============================================================
\documentclass[12pt, a4paper]{article}

% ── Encoding y fuentes ────────────────────────────────────────
\usepackage[utf8]{inputenc}
\usepackage[T1]{fontenc}
\usepackage[spanish]{babel}

% ── Matemática y física ───────────────────────────────────────
\usepackage{amsmath, amssymb}
\usepackage{siunitx}           % unidades del SI: \SI{9.8}{m/s^2}
\sisetup{
    output-decimal-marker = {,},
    separate-uncertainty  = true,
}

% ── Tablas y figuras ──────────────────────────────────────────
\usepackage{booktabs}          % \toprule, \midrule, \bottomrule
\usepackage{graphicx}
\usepackage{float}             % opción [H] para figuras
\usepackage{caption}

% ── Bibliografía ──────────────────────────────────────────────
\usepackage[backend=biber, style=numeric, sorting=none]{biblatex}
\addbibresource{referencias.bib}

% ── Otros ─────────────────────────────────────────────────────
\usepackage{hyperref}
\usepackage{geometry}
\geometry{margin=2.5cm}

% ==============================================================
\title{Práctica NN: Título del experimento}
\author{Nombre Apellido}
\date{\today}
% ==============================================================

\begin{document}

\maketitle

% ── Resumen ───────────────────────────────────────────────────
\begin{abstract}
Breve descripción del experimento, metodología principal y resultado central.
\end{abstract}

% ── 1. Introducción ───────────────────────────────────────────
\section{Introducción}

Contexto teórico mínimo necesario para entender el experimento.
Citar referencias relevantes \cite{ejemplo}.

% ── 2. Montaje experimental ───────────────────────────────────
\section{Montaje experimental}

Describir el setup, instrumentos y condiciones de medición.

% ── 3. Análisis de datos ──────────────────────────────────────
\section{Análisis de datos}

\subsection{Procesamiento}

Describir cómo se procesaron los datos (filtros, ajustes, etc.).

\subsection{Propagación de incertezas}

Expresar las incertezas de las magnitudes derivadas.

% Ejemplo de figura
% NOTA: compilar desde la carpeta informe/ con: pdflatex informe.tex
%       La ruta ../figuras/ asume que el directorio de trabajo es informe/
\begin{figure}[H]
    \centering
    \includegraphics[width=0.7\textwidth]{../figuras/fig_ejemplo.pdf}
    \caption{Descripción de la figura.}
    \label{fig:ejemplo}
\end{figure}

% ── 4. Resultados ─────────────────────────────────────────────
\section{Resultados}

Presentar los resultados principales con sus incertezas.
Por ejemplo: $g = \SI{9.79 \pm 0.05}{m/s^2}$.

% Ejemplo de tabla
\begin{table}[H]
    \centering
    \caption{Resultados del ajuste.}
    \label{tab:resultados}
    \begin{tabular}{lcc}
        \toprule
        Magnitud & Valor & Incerteza \\
        \midrule
        $g$ & \SI{9.79}{m/s^2} & \SI{0.05}{m/s^2} \\
        \bottomrule
    \end{tabular}
\end{table}

% ── 5. Conclusiones ───────────────────────────────────────────
\section{Conclusiones}

Discutir los resultados en relación con los valores esperados o teóricos.
Mencionar fuentes de error sistemático y posibles mejoras.

% ── Bibliografía ──────────────────────────────────────────────
\printbibliography

\end{document}
